\documentclass[11pt,a4paper]{article}

\usepackage[utf8]{inputenc}
\usepackage[T1]{fontenc}
\usepackage[english]{babel}
\usepackage{ctex}
\usepackage[margin=2.5cm]{geometry}
\usepackage{parskip}
\usepackage{amsmath}
\usepackage{amssymb}
\usepackage{bm}
\usepackage[numbers,sort&compress]{natbib}
\renewcommand*{\bibfont}{\footnotesize}
\usepackage{xcolor}
\usepackage{hyperref}
\hypersetup{
  colorlinks=true,
  linkcolor=blue,
  urlcolor=cyan,
  pdftitle={机器灌注条件下细胞外囊泡调控移植器官缺血再灌注损伤及早期排斥风险：机制、标志物与转化应用},
  pdfsubject={Systematic Literature Review},
}

\title{机器灌注条件下细胞外囊泡调控移植器官缺血再灌注损伤及早期排斥风险：机制、标志物与转化应用}
\author{}
\date{2026-06-03}

\providecommand{\tightlist}{%
  \setlength{\itemsep}{0pt}\setlength{\parskip}{0pt}}

\begin{document}
\maketitle
\tableofcontents
\newpage

\begin{abstract}
缺血再灌注损伤（ischemia-reperfusion injury,
IRI）是实体器官移植从供体获取、保存、植入到早期免疫激活之间的共同病理枢纽。近五年，低温含氧机器灌注、常温机器灌注与控制性复温逐步把器官保存从静态冷保存推进为可实时监测和可干预的离体平台；与此同时，细胞外囊泡（extracellular
vesicles,
EVs）被重新定位为损伤信号、免疫调控载体和潜在治疗递送系统。现有证据提示，EV
cargo 中的
miRNA、蛋白、脂质和线粒体相关成分能够影响内皮糖萼、补体、炎症小体、巨噬细胞极化和细胞死亡；机器灌注则为
EV
的器官靶向递送、灌注液液体活检和短时药效评价提供了独特窗口。真正的转化挑战在于
EV
来源、剂量、纯度、效价、器官内分布、长期免疫后果和监管质量标准仍未统一。未来
3-5 年最有价值的方向，是把``EV
治疗''从泛化的抗炎叙事推进到灌注平台上的可计量药代、可解释标志物和小样本自适应临床研究。
\end{abstract}

\section{引言}\label{ux5f15ux8a00}

IRI
并不是单纯的氧化应激事件，而是冷缺血、复温、线粒体损伤、内皮屏障破坏、补体启动、固有免疫细胞募集和后续适应性免疫放大的连续过程。肝、肾、心、肺等器官在解剖和代谢上不同，但早期损伤均可通过
DAMPs、线粒体 DNA、内皮糖萼脱落、NLRP3
炎症小体和微循环障碍汇入相似的免疫激活网络\cite{mechanisticinsightand2023, rationalefortargeting2025, theendothelialglycocalyx2021}。临床上，这一网络会转化为移植物功能延迟、原发无功能、胆道并发症、急性排斥易感性以及慢性移植物损伤的底层土壤。

机器灌注的重要性在于它改变了干预时点。静态冷保存让医生主要在植入后处理损伤，而机器灌注允许在移植前评估代谢恢复、微循环阻力、胆汁/尿液生成、乳酸清除、酶学释放和灌注液分子标志物\cite{exvivokidney2021, currentevidenceand2023, perfusiontechniquesin2024}。这使器官保存从``运输技术''变成``短时体外药理学平台''。EV
的兴起恰好与这一平台相遇：EV
可以是损伤读数，也可以是治疗载体；可以来源于供体器官、免疫细胞、间充质
stromal cells，也可以被工程化装载抗炎或促修复
cargo\cite{extracellularvesiclesas2023, organrepairand2023}。

\section{IRI
的免疫机制：从线粒体失衡到排斥风险}\label{iri-ux7684ux514dux75abux673aux5236ux4eceux7ebfux7c92ux4f53ux5931ux8861ux5230ux6392ux65a5ux98ceux9669}

IRI 的核心机制是能量失败后再供氧所触发的非线性放大。冷缺血阶段 ATP
耗竭、离子泵失衡和线粒体电子传递链损伤已经埋下伏笔；复温和再氧合阶段，ROS
暴发、钙超载、膜通透性转换孔开放和线粒体 DAMPs 释放共同激活内皮、Kupffer
细胞、肾小管细胞和组织驻留免疫细胞\cite{mitochondrialmetabolismand2024, focusingonischemic2024}。补体是这一过程中特别值得靶向的层级，因为它既连接先天免疫，也能通过
C3a/C5a、膜攻击复合物和凝血串扰放大微血管损伤\cite{rationalefortargeting2025}。

EV 在这一网络中具有双重身份。损伤细胞释放的 EV
可携带氧化脂质、miRNA、线粒体片段和组织特异蛋白，成为``损伤正在发生''的信号；治疗性
EV 则可能通过转运抗凋亡
miRNA、免疫调节蛋白和线粒体稳态相关成分，削弱炎症与细胞死亡。问题在于这两类
EV 在粒径、来源、cargo
和生物学效应上并不天然可分。若没有来源追踪和效价定义，EV
研究容易把``检测到囊泡变化''误读为``囊泡驱动病理''。

\section{机器灌注作为 EV
递送窗口}\label{ux673aux5668ux704cux6ce8ux4f5cux4e3a-ev-ux9012ux9001ux7a97ux53e3}

机器灌注的最大优势是把治疗递送限制在器官层面。对于
MSC-EV、尿液祖细胞来源 EV
或工程化纳米囊泡，离体灌注可以减少全身给药的稀释、免疫副作用和非靶器官分布，同时允许在植入前洗脱未结合成分\cite{applicationofmesenchymal2021, extracellularvesiclesderived2022, dynamicconditioningof2024}。在扩展标准供肾和
DCD 心脏模型中，MSC-EV 通过 HOPE
或相关灌注策略递送后显示出抗炎、抗凋亡和改善微循环的信号；猪 DCD
心脏模型中的初步研究则支持 EV 对心肌 IRI
具有保护潜力\cite{protectiveeffectsof2024}。

但当前证据仍以小动物、大动物和离体概念验证为主。EV
在灌注液中的稳定性、血管内皮摄取、实质细胞渗透深度、胆道/肾小管分布，以及不同温度、氧合状态和剪切力下的
cargo
保留，仍缺少可比较的数据。常温机器灌注更接近生理代谢，适合评价即时药效和组织摄取；低温含氧灌注对线粒体恢复和
ROS 预处理更有优势，但 EV 内吞和细胞代谢反应可能受低温限制。因此，EV
递送不能脱离灌注模式讨论。

\section{灌注液 EV
与分子标志物}\label{ux704cux6ce8ux6db2-ev-ux4e0eux5206ux5b50ux6807ux5fd7ux7269}

EV
的另一条转化路径是标志物。灌注液比外周血更接近供体器官，背景噪音低，可连续采样，适合捕捉内皮、胆管、肾小管、心肌和免疫细胞释放的时间序列信号。近年研究开始把
EV、miRNA、代谢物、蛋白组和细胞游离 DNA
共同纳入器官质量评价框架\cite{alectinaffinity2025, abiomarkerframework2024, massspectrometrybased2023}。灌注液
EV/miRNA
被清除或富集后的变化，也提示``过滤有害囊泡''可能与``递送有益囊泡''形成互补策略\cite{thehemopurifierremoves2024}。

真正可用的标志物不应只是单个 miRNA
或单个炎症因子，而应服务三个临床问题：器官是否适合移植、是否需要灌注期修复、植入后早期是否需要强化监测。为此，灌注液
EV
指标需要与灌注动力学、组织病理、术后功能延迟、dd-cfDNA、供受体免疫风险和影像/病理端点耦合。仅靠相关性不足以改变临床决策；需要前瞻性阈值、阴性预测价值和干预后风险重分类证据。

\section{讨论：主要瓶颈}\label{ux8ba8ux8bbaux4e3bux8981ux74f6ux9888}

第一，EV 产品本身高度异质。MSC
来源、培养条件、缺氧预处理、纯化方法、粒径分布、蛋白污染、RNA cargo
和效价检测都会影响结果。第二，剂量单位不统一，粒子数、蛋白量、细胞当量和功能效价之间缺少换算。第三，安全性不仅包括急性毒性，还包括促凝、促血管生成、免疫耐受过度、潜在促肿瘤和感染载体风险。第四，机器灌注研究常常样本量小、器官来源不均一、灌注参数差异大，导致跨研究比较困难。

第五，也是最容易被低估的一点：EV
可能同时是治疗物和生物标志物。如果治疗性 EV 改变了灌注液 EV
谱，那么后续以 EV
作为质量读数时必须区分内源性释放和外源性给药残留。未来研究需要把可追踪标签、单囊泡分析和组织空间组学纳入方案，否则很难证明治疗性
EV 到达了正确的细胞并产生了正确的效应。

\section{展望：未来 3-5
年研究方向}\label{ux5c55ux671bux672aux6765-3-5-ux5e74ux7814ux7a76ux65b9ux5411}

\begin{enumerate}
\def\labelenumi{\arabic{enumi}.}
\item
  建立``灌注液 EV 药代学''。在 NMP、HOPE 和控制性复温中系统测定 EV
  的半衰期、内皮摄取、组织分布、洗脱效率和 cargo
  保留，以器官为单位建立剂量-暴露-效应曲线。
\item
  发展可追踪、可放行的 EV 质控体系。建议把粒径/浓度、污染蛋白、关键
  miRNA/蛋白
  cargo、促凝风险、内毒素、无菌、效价模型和批间一致性作为最小放行集。
\item
  将 EV 递送与补体/线粒体靶向联合。单一 EV 抗炎效应可能不足以压制 IRI
  网络；更现实的策略是在灌注期联合
  EV、补体抑制、线粒体保护和内皮糖萼保护，并用灌注液读数实时判断是否达到生物学效应\cite{rationalefortargeting2025, mesenchymalstemcells2024}。
\item
  构建``灌注液-组织-外周血''三层预测模型。灌注液
  EV/miRNA、代谢组和蛋白组负责移植前风险；植入后 dd-cfDNA
  和炎症指标负责早期损伤；病理和长期功能负责模型校准。
\item
  设计小样本自适应临床前-临床桥接研究。高风险 DCD
  或扩展标准供体器官最适合作为早期场景，端点应优先选择安全性、器官利用率、功能延迟、早期
  dd-cfDNA 和组织损伤，而不是过早追求长期生存差异。
\end{enumerate}

\section{结论}\label{ux7ed3ux8bba}

EV 与机器灌注的交汇，让移植 IRI
研究获得了一个可操作的转化场景。机器灌注提供时间、空间和监测窗口；EV
提供可塑的信号载体和治疗载体。短期内，灌注液 EV/miRNA
作为器官质量读数可能比 EV 治疗更快进入临床决策；中期内，标准化 EV
产品与灌注期联合干预有望成为边缘供体器官修复的重要方向。该领域未来的成败，将取决于能否把漂亮的抗炎叙事变成可计量、可重复、可监管、可改变移植决策的证据链。

\nocite{*}
\bibliographystyle{gbt7714-numerical}
\bibliography{ev-machine-perfusion-transplant-iri_参考文献}

\end{document}
